\documentclass[11pt]{article}

\usepackage[margin=1in]{geometry}
\usepackage{amsmath}
\usepackage{amssymb}
\usepackage{booktabs}
\usepackage{hyperref}

\title{Executive Summary: Dupire vs. Neural-Network Delta Hedging}
\author{Oscar Arevalo, Micheal O Cobhthaigh, Alec Hewitt, Ziqing Zhang}
\date{\today}

\begin{document}

\maketitle

\section*{Project Objective}

This project compares two approaches to delta hedging a short SPY call option:
a classical Dupire local-volatility approach and a neural-network-based hedge
ratio model. The goal is not to build a live trading system, but to evaluate how
two different modeling frameworks translate market data into hedge ratios and
how those hedge ratios behave in an out-of-sample daily rebalancing simulation.

The central question is:
\begin{quote}
How well can a classical local-volatility model and a data-driven neural network
control the risk of a short option position through dynamic stock hedging?
\end{quote}

\section*{Data Pipeline}

The data construction is handled in \texttt{gen\_dat.ipynb}. The notebook reads
SPY option archive files from \texttt{archive/*.parquet}, cleans the raw quote
columns, and converts the option records into a project-standard schema. The
analysis uses call options with fields such as quote date, expiration, strike,
bid, ask, last price, implied volatility, volume, and option Greeks from the
archive.

The generated datasets are then consumed by \texttt{dupire\_vs\_nn\_hedging.ipynb}.
That notebook requires \texttt{data/historical\_options.csv}; if this file is not
available, the notebook stops rather than silently using synthetic option data.
The current run uses real historical option data generated from the SPY archive.

The loader applies consistent cleaning:
\begin{itemize}
    \item restricts the study to call options;
    \item computes mid price as $(\text{bid}+\text{ask})/2$;
    \item computes time to maturity $T$ in years;
    \item computes moneyness $K/S$ and log-moneyness $\log(K/S)$;
    \item filters to positive bid/ask quotes, positive open interest, positive
    implied volatility, maturities between 7 and 365 days, and moneyness between
    0.8 and 1.2.
\end{itemize}

The stock data are SPY closing prices with daily returns, log returns, and
20-day and 60-day realized volatility estimates. These stock features are merged
onto the option observations by date.

\section*{Financial Motivation}

The Black--Scholes model assumes constant volatility:
\[
    dS_t = rS_t\,dt + \sigma S_t\,dW_t.
\]
In observed option markets, however, the volatility implied by option prices
varies across strike and maturity:
\[
    \sigma_{\mathrm{imp}} = \sigma_{\mathrm{imp}}(K,T).
\]
This implied volatility surface shows that one constant volatility parameter is
too restrictive for fitting the cross-section of traded option prices.

The presentation in \texttt{dupire\_presentation/} motivates the Dupire model
from this observation. Dupire converts an option price or implied-volatility
surface into a local-volatility function that depends on stock price and time:
\[
    dS_t = rS_t\,dt + \sigma_{\mathrm{loc}}(S_t,t)S_t\,dW_t.
\]
This gives a dynamic stock-price model calibrated to the observed option surface
under ideal smooth and arbitrage-free conditions.

\section*{Dupire Local-Volatility Method}

The Dupire method is implemented in \texttt{src/dupire.py}. The project first
fits an implied-volatility interpolator over log-moneyness and maturity:
\[
    x = \log(K/S), \qquad \sigma_{\mathrm{imp}} = \sigma_{\mathrm{imp}}(x,T).
\]
The code uses linear interpolation where possible and nearest-neighbor fallback
where linear interpolation is unavailable. Implied volatilities are clipped to a
practical range to avoid extreme numerical behavior.

Next, the code builds a grid of log-moneyness and maturity values, converts
log-moneyness back to strikes using
\[
    K = S_0 e^x,
\]
and prices calls on the grid using the Black--Scholes formula with the
interpolated implied volatility:
\[
    C(K,T) =
    S_0 N(d_1) - K e^{-rT}N(d_2).
\]
This produces a numerical call price surface $C(K,T)$.

The local variance is then estimated with the Dupire formula:
\[
    \sigma_{\mathrm{loc}}^2(K,T)
    =
    \frac{
    \frac{\partial C}{\partial T}
    + rK\frac{\partial C}{\partial K}
    }{
    \frac{1}{2}K^2\frac{\partial^2 C}{\partial K^2}
    }.
\]
The derivatives are computed numerically using finite differences. When the
formula produces invalid or negative local variance, the implementation falls
back to implied variance and clips the resulting local volatility to a stable
range.

\subsection*{Dupire Delta Proxy}

A fully model-consistent Dupire delta would solve the local-volatility pricing
problem for the option value $V(S,t)$ and then compute
\[
    \Delta_{\mathrm{Dupire}} = \frac{\partial V}{\partial S}.
\]
For speed and interpretability, this project uses a practical proxy: it looks up
the local volatility at the option's moneyness and maturity, then plugs that
volatility into the Black--Scholes call delta formula:
\[
    \Delta_{\mathrm{proxy}}
    =
    N\!\left(
    \frac{
    \log(S/K) + \left(r+\frac{1}{2}\sigma_{\mathrm{loc}}^2\right)T
    }{
    \sigma_{\mathrm{loc}}\sqrt{T}
    }
    \right).
\]
This is an educational approximation rather than a production local-volatility
delta engine. As a check in \texttt{Test/Dupire/Dupire\_test.ipynb}, a
finite-difference Monte Carlo estimate under local-volatility dynamics produced
a numerical delta of approximately $0.47$ versus a proxy value of approximately
$0.54$ for an at-the-money 30-day example. The values are directionally similar,
but not identical, which is expected because the proxy is not the full
local-volatility PDE delta.

\section*{Neural-Network Hedge Method}

The neural-network method is implemented in \texttt{src/nn\_hedge.py} and
\texttt{src/features.py}. The neural network is a feedforward regression model
with three hidden layers:
\[
    64 \rightarrow 32 \rightarrow 16
\]
using ReLU activations and a single scalar output. The output is interpreted as
a predicted option hedge ratio.

The model uses the following feature set:
\begin{itemize}
    \item moneyness and log-moneyness;
    \item time to maturity $T$;
    \item implied volatility;
    \item 20-day and 60-day realized stock volatility;
    \item Black--Scholes call delta;
    \item option mid price;
    \item relative bid-ask spread;
    \item volume and open interest;
    \item recent stock return and absolute recent stock return.
\end{itemize}

The Black--Scholes delta is included as one input feature, not as the main target
in the historical run. The target is a realized next-period delta estimate:
\[
    \Delta_{\mathrm{realized}}
    \approx
    \frac{\Delta C}{\Delta S}
    =
    \frac{C_{t+1}-C_t}{S_{t+1}-S_t},
\]
computed within repeated observations of the same strike and expiration. The
target is clipped to a practical range of $[-2,2]$. If repeated contract
observations are unavailable, the code warns and falls back to Black--Scholes
delta as a proxy target.

The training procedure sorts observations by date, uses the earliest 70\% for
training, the next 15\% for validation, and reserves the remaining period for
out-of-sample testing. Features are standardized with \texttt{StandardScaler}.
Training uses mean squared error loss, Adam optimization, and early stopping
with patience 35.

\section*{Backtest Design}

Both models are evaluated on the same out-of-sample short-call portfolio. The
portfolio selection routine chooses a consistently observed call contract and
prefers contracts close to at-the-money. The default project configuration is:
\begin{itemize}
    \item underlying: SPY;
    \item risk-free rate: $r=4\%$;
    \item position: one short call contract;
    \item contract multiplier: 100;
    \item rebalancing: daily;
    \item transaction cost: 1 basis point on stock trades;
    \item initial cash: 0.
\end{itemize}

For each day in the out-of-sample portfolio, the strategy converts the option
delta into the number of SPY shares needed to offset the short option's delta:
\[
    \text{shares} =
    -(\text{option quantity}) \times \Delta \times \text{contract multiplier}.
\]
The backtest tracks cash, stock position, option value, portfolio value, daily
P\&L, hedging error, transaction costs, and turnover.

The portfolio value is the marked value of the combined hedged position:
\[
    \text{portfolio value}
    =
    \text{short option value}
    +
    \text{stock hedge value}
    +
    \text{cash}.
\]
This should not be interpreted as directly comparable to a long-only SPY
buy-and-hold portfolio, because the hedging strategies include a short option
liability and dynamic stock rebalancing.

\section*{Results Summary}

The current run evaluates 63 out-of-sample portfolio observations. The summary
metrics saved in \texttt{results/strategy\_comparison.csv} are:

\begin{center}
\begin{tabular}{lrrrr}
\toprule
Model & Total P\&L & Sharpe & Mean Abs. Error & Costs \\
\midrule
Dupire Delta Hedge & 291.24 & 1.66 & 34.61 & 8.17 \\
Neural-Network Delta Hedge & 297.33 & 0.74 & 74.18 & 10.90 \\
\bottomrule
\end{tabular}
\end{center}

In this run, the neural-network hedge has slightly higher total P\&L, but the
Dupire hedge has materially lower P\&L volatility, lower mean absolute hedging
error, lower RMSE hedging error, lower maximum drawdown, lower transaction
costs, and lower turnover. The Dupire strategy therefore appears more stable as
a hedge in this particular experiment, while the neural network produces a
somewhat higher total P\&L at the cost of noisier hedge behavior.

The notebook also includes an optional SPY buy-and-hold market benchmark. This
benchmark is useful only as market context. It is not a competing hedging
strategy because it does not include a short option liability, does not compute a
hedge ratio, and has no option hedging error.

\section*{Key Takeaways}

\begin{itemize}
    \item The project demonstrates an end-to-end comparison between a classical
    financial model and a learning-based hedge ratio model.
    \item The Dupire method is mathematically grounded in the implied-volatility
    surface, but it is sensitive to noisy market data and numerical
    differentiation.
    \item The neural network can learn hedge ratios from market features and
    realized option-stock movements, but it requires careful validation and may
    produce less stable hedge behavior.
    \item Hedging quality should be judged with risk and error metrics, not only
    total P\&L. Lower hedging error, lower turnover, lower transaction costs, and
    lower drawdown are central to evaluating hedge effectiveness.
    \item Better hedging does not necessarily imply a profitable trading
    strategy. The project studies hedge performance, not an alpha signal or
    standalone trading strategy.
\end{itemize}

\section*{Limitations and Future Work}

Several limitations are important for interpreting the results. The Dupire
surface is an educational approximation and uses finite differences on an
interpolated surface, which can amplify quote noise. The delta used in the
backtest is a Black--Scholes-with-local-volatility proxy rather than a full
PDE-derived local-volatility delta. The neural network is relatively simple and
is trained on realized next-period delta estimates, which can themselves be
noisy. Finally, the backtest uses one selected short-call portfolio and should
not be interpreted as a comprehensive live trading evaluation.

Natural extensions include stronger surface smoothing, a full PDE-based Dupire
delta calculation, broader option portfolio construction, richer transaction
cost modeling, cross-validation across market regimes, and comparison with
additional hedge models.

\end{document}
